% Exploratory conditional reduction, outside the frozen CP27 paper.
% Root written after its 23:00 UTC source and local-proof checks.
\subsection{A conditional single-orbit inertia route for odd degrees}

This note does not assert a new monodromy image or a smaller proved main
gate. The required global branch-divisor assertions below are unproved.
The two elementary lemmas have complete direct proofs and an independent
proof audit. In particular an irreducible branch divisor alone is insufficient.

\begin{lemma}[One conjugacy class of transvections]
\label{lem:root-single-transvection-class}
\status{Proved; complete direct proof and independent audit passed.}
Let $V$ be a nonzero finite-dimensional vector space over a field $k$, and
let $G\subseteq\operatorname{GL}_k(V)$ be generated by a single
$G$-conjugacy class $\mathcal C$ of transvections. If
$V^G=0=(V^\vee)^G$, then $V$ is an irreducible $k[G]$-module.

\end{lemma}
\begin{proof}
Suppose $0\ne S\subsetneq V$ is $G$-stable. If every element of
$\mathcal C$ acts identically on $S$, its generators do so, hence
$S\subseteq V^G=0$, a contradiction. Otherwise choose $t\in\mathcal C$
and $s\in S$ with $(t-1)s\ne0$. Since $t-1$ has rank one,
$\operatorname{im}(t-1)=k(t-1)s\subseteq S$. Conjugacy and stability
then give $\operatorname{im}(u-1)\subseteq S$ for every
$u\in\mathcal C$. Thus every generator acts identically on $V/S$.
Any nonzero linear functional on $V/S$ pulls back to a nonzero element
of $(V^\vee)^G$, again a contradiction. This proves irreducibility.

\end{proof}
\begin{lemma}[Local corank-one Lang inertia]
\label{lem:root-corank-one-Lang-inertia}
\status{Proved; complete direct proof and independent audit passed.}
Let $R$ be a strictly henselian discrete valuation ring of characteristic
two, with fraction field $K$, and let $E$ be a finite free $R$-module
with a Frobenius-semilinear operator $F$. Suppose the solution spaces
of $F(v)=v$ on a geometric generic and special fibre have respective
$\mathbf F_2$-dimensions $r$ and $r-1$. Then the inertia image on the
generic solution space $V$ fixes pointwise a hyperplane $W$ and embeds
in the additive group
$\operatorname{Hom}_{\mathbf F_2}(V/W,W)$. Every nonidentity inertia
element is a transvection. Its image need not be cyclic.

\end{lemma}
\begin{proof}
In a basis, the Lang zero scheme has equations $A x^{[2]}-x=0$.
Their Jacobian is $-I$, so it is etale over $R$, even though it need
not be finite over $R$. Each special solution therefore has a unique
lift through its henselian neighbourhood to an $R$-section. A relation
between the generic lifts reduces to the same relation in the special
fibre; hence the lifts form an $(r-1)$-dimensional subspace $W$.
They are $K$-rational, so inertia fixes $W$ pointwise. On the
one-dimensional quotient $V/W$, inertia acts through
$\mathbf F_2^\times=\{1\}$. Consequently every $g-1$ factors from
$V/W$ into $W$, and these factorizations multiply by addition since
$(g-1)(h-1)=0$. A nonzero factorization has rank one and square zero.
This proves the claims. In particular the permitted image can have
order as large as $2^{r-1}$, and its nonzero elements can have different
image lines inside $W$.

\end{proof}

\par\smallskip\noindent\textbf{Precise unresolved geometric step.}
For a smooth simply connected compactification of the parameter space,
normal generation by divisorial inertia follows only after proving
that the finite quotient cover has no remaining divisorial branch and
applying purity. If the entire branch support were one irreducible
divisor, and its generic inertia were cyclic of order two, this would
give the single class required in the first lemma. More generally one
needs normal generation and conjugacy of all generating transvections,
not merely conjugacy of the generic inertia subgroups.

Neither the required branch-support statement nor cyclicity/conjugacy
has been proved for the remaining odd fixed-surface systems.
The monodromy reviewer reports that the degree-one relative determinant
already has the form $q(s)^2Q(s)$ with two cubic components, and that its
moving-kernel inertia is an elementary abelian group rather than a
single cyclic generator. The previously withdrawn higher-degree
alternating/Pfaffian model supplies no permissible higher-degree
factorization. Thus the existing odd-plus-even parabolic countermodel
is not excluded by the presently proved geometry.

The opened primary comparison was Brumer and Kramer~\cite{BrumerKramer2018LargeImage}, arXiv:1701.01890v1,
Theorem 2.1.1 and its reduction on printed p.3, and Proposition 2.3.4
with complete proof on printed p.5. Those are symplectic and arithmetic
statements; none is transferred to the present characteristic-two
unit-root family. The elementary criterion above is proved directly.

% Post-freeze exploratory route obstruction, not a new full-monodromy theorem.
\subsection{Marked normalization can make triple-collision inertia trivial}

\begin{proposition}[Trivial inertia with marked constant-rank normalization]
\label{prop:root-marked-normalization-inertia}
\status{Proved; complete direct proof and independent audit passed.}
Let $R$ be a strictly henselian discrete valuation ring of characteristic
two, with fraction field $K$. Let $C\to\operatorname{Spec}R$ be a smooth
proper geometrically connected curve whose geometric generic and special
$2$-ranks are the same integer $f$. Let $Y/K$ be a proper geometrically
connected nodal curve with normalization $C_K$. Assume its $\delta$ nodes
are $K$-rational and that the two branches of each node are given by
$R$-sections of $C$, all distinct on the generic fibre. The sections may
coalesce on the special fibre. Then the inertia action on
$H^1_{\mathrm{et}}(Y_{\overline K},\mathbf F_2)$ is trivial.

Here the claim concerns this marked-normalization situation only. It does
not assert that an arbitrary triple-point family has this form.

\end{proposition}
\begin{proof}
The $R$-module $E=H^1(C,\mathcal O_C)$ is finite free and commutes with
base change. Its Frobenius-semilinear operator has an etale Lang zero
scheme, since the equations $A x^{[2]}-x=0$ have Jacobian $-I$.
Every special fixed point lifts uniquely through its henselian
neighbourhood to an $R$-point. The special and generic solution spaces
both have $2^f$ elements, so these lifts exhaust all geometric generic
solutions. Thus every generic Frobenius-fixed cohomology class is defined
over $R$.

The map $F-1:R\to R$ is surjective by the strict henselian property
applied to $z^2-z-a$. Since $H^0(C,\mathcal O_C)=R$, the Artin--Schreier
exact sequence identifies $H^1_{\mathrm{et}}(C,\mathbf F_2)$ with
$\ker(F-1:E\to E)$. Consequently every geometric generic
$\mathbf F_2$-torsor on $C_K$ is the restriction of a torsor on $C$.
Its pullback to each marked section is a finite etale torsor over the
strictly henselian ring $R$, hence is trivial. Choose these
trivializations separately at all marked branches.

An equivariant identification of the two torsor fibres above a generic
node has exactly two choices. Both choices are defined over $K$ by the
chosen trivializations. Any collection of these choices glues the
normalization torsor to a finite etale $\mathbf F_2$-torsor on $Y$.
For completeness, its finite algebra is the subalgebra of the finite
pushforward of the normalization-cover algebra for which paired residue
values agree under the chosen identifications. At a node, after strict
henselization, the base algebra is the fibre product of the two branch
algebras over the residue field. The chosen matching makes the cover
algebra the product of two copies of this fibre product. It is therefore
finite etale at the node, as well as away from the nodes. Conversely a
torsor on $Y$ recovers exactly its normalization torsor and these
equivariant identifications, so every geometric torsor is obtained.

All geometric cohomology classes on $Y$ thus have representatives defined
over $K$. Inertia fixes every class, proving the assertion. The argument
does not require the chosen generic pairwise identifications to define a
torsor on a possible special curve $Y_0$.

\end{proof}

\par\smallskip\noindent\textbf{Effect on the proposed triple-packet route.}
When three generic nodal identifications merge into an ordinary triple
point, a new cycle-consistency condition can exclude some collections of
signs from extending across the special point. Under the hypotheses
above, all those signs were nevertheless already defined over $K$.
Their failure to extend is therefore compatible with trivial inertia:
it is not evidence of a nonzero translation in monodromy. The proposed
pure-translation triple-packet argument cannot be used without proving
an additional phenomenon outside this marked constant-rank normalization
situation.

Root's gluing objection was communicated before 23:04 UTC on
19 November 2026, and the monodromy reviewer independently agreed with
the argument in the same interval. This fully written note was retained
before the checkpoint27 archive; it was outside that frozen paper and
outcome table and is now included with its exact limited hypotheses. No new numerical computation or external theorem read is
claimed. It records a failed research route, not main-gate movement.

% A limitation of a proposed route, not a geometric counterexample to A.4.
\subsection{Rational algebraic monodromy does not determine the residual image}

\begin{proposition}[A compact open residual countermodel]
\label{prop:root-rational-residual-obstruction}
\status{Proved; complete direct proof and independent audit passed.}
Let \(p\) be a prime, \(d\ge1\), and
\(P\subseteq\operatorname{GL}_d(\mathbf F_p)\) any subgroup.
Write \(L=\mathbf Z_p^d\), and let
\[
 \operatorname{red}:\operatorname{GL}_{\mathbf Z_p}(L)
       \longrightarrow\operatorname{GL}_{\mathbf F_p}(L/pL)
\]
be reduction, a continuous surjective group homomorphism. Then
\(\Gamma=\operatorname{red}^{-1}(P)\), with its natural profinite
topology, acts continuously on \(L\), has residual image exactly
\(P\), and has Zariski closure
\(\operatorname{GL}_{d,\mathbf Q_p}\) on
\(L\otimes_{\mathbf Z_p}\mathbf Q_p\).
In particular its rational representation is absolutely irreducible,
irrespective of reducibility of the residual representation.
\end{proposition}
\begin{proof}
An invertible matrix over \(\mathbf F_p\) admits an entrywise lift
whose determinant is a \(p\)-adic unit; this proves surjectivity of
\(\operatorname{red}\), hence the exact residual image. The inverse
image of a subgroup of a finite discrete group is a compact open
subgroup. It contains \(1+p\operatorname{Mat}_d(\mathbf Z_p)\).

Every polynomial over \(\mathbf Q_p\) vanishing on the latter
set is zero. Indeed substitute \(x_{ij}=\delta_{ij}+p y_{ij}\).
The resulting polynomial vanishes on the Cartesian product
\(\mathbf Z_p^{d^2}\). In one variable infinitely many roots force
it to vanish identically; induction on the variables proves the same
assertion for the product. The substitution is an invertible affine
change over \(\mathbf Q_p\). A regular function on
\(\operatorname{GL}_{d,\mathbf Q_p}\) is a polynomial divided by
a power of the determinant, so the same reasoning proves the stated
Zariski density. The argument remains true after extending scalars to
an algebraic closure: expand the finitely many coefficients in a
\(\mathbf Q_p\)-basis of their span and apply the preceding argument
to each coefficient polynomial. Any invariant subspace in the
algebraic closure is consequently invariant under the full general
linear group, whose natural representation is irreducible.
\end{proof}

\begin{example}[The absence of fixed vectors and covectors still permits a parabolic]
\label{ex:root-rational-residual-parabolic}
\status{Proved; direct proof, full independent audit and exact finite supplement passed.}
For \(r>s\ge2\), put \(d=r+s\) and
\[
 P=\left\{
  \begin{pmatrix} A&B\\0&D\end{pmatrix}:
  A\in\operatorname{GL}_r(\mathbf F_p),\quad
  D\in\operatorname{GL}_s(\mathbf F_p),\quad
  B\in\operatorname{Mat}_{r\times s}(\mathbf F_p)
 \right\}.
\]
This is a proper subgroup preserving the proper nonzero subspace
\(U=\mathbf F_p^r\oplus0\). For the natural module
\(V=\mathbf F_p^d\), nevertheless
\(V^P=0=(V^\vee)^P\), and \(P\) contains the full natural
\(\operatorname{GL}_r(\mathbf F_p)\) block with
\(r>d/2\). Its inverse image in
\(\operatorname{GL}_d(\mathbf Z_p)\) has the full rational
algebraic monodromy of the preceding proposition.
\end{example}
\begin{proof}
Block multiplication and inversion preserve the displayed form, so
it is a subgroup. An elementary matrix carrying one vector of \(U\)
to a vector with nonzero second component does not preserve \(U\),
proving properness. The subgroup of block diagonal matrices acts
independently on the two summands. For any \(m\ge2\), the natural
\(\operatorname{GL}_m(\mathbf F_p)\)-module has no fixed nonzero
vector: a nonzero coordinate can be added to a different coordinate
by an elementary transvection. The same argument after transpose
applies to covectors. Apply this separately to both diagonal blocks.
The large natural block is the subgroup with \(D=1,B=0\), and
the final assertion follows from
Proposition~\ref{prop:root-rational-residual-obstruction}.

The retained program
\path{code/check_rational_residual_parabolic.py}
checks the explicit finite instance \(p=2,r=3,s=2\). It enumerates
the full displayed subgroup, checks every member's rank and invariant
subspace, and checks the fixed vectors and covectors against elementary
generators. Its output states exactly those finite claims. The
infinite \(p\)-adic density and all-rank assertions are proved above.
\end{proof}

\par\smallskip\noindent\textbf{Scope of the literature route.}
We examined the original preprint of D'Addezio
\cite[Theorem~4.4.12, printed pp.~23--24;
Proposition~5.1.4 and Theorem~5.1.6, printed p.~25]
{DAddezio2023Parabolicity}. The final proofs at these locators were
read, but their deeper dependencies were not all independently
audited. The statements concern rational isocrystals and rational
algebraic monodromy; Theorem~5.1.6 in particular states rational
semisimplicity for an abelian scheme of constant Newton slopes.
We do not import these results as a residual-image theorem. Even
full rational algebraic monodromy and absence of residual fixed
vectors and covectors leave the explicit parabolic above possible.
Additional integral or residual geometry is necessary for A.4.
This example is an abstract continuous representation; it does not
assert its realization by the hypersurface family in the question,
and it does not shrink the remaining geometric main gate.

